\section{A Publication Agenda Beyond Facelift Revision}

\begin{table*}[t]
\centering
\footnotesize
\begin{tabular}{L{0.21\linewidth} L{0.20\linewidth} L{0.20\linewidth} L{0.27\linewidth}}
\toprule
Research thrust & Near-term technical question & Why biomolecular AI is a strong fit & Failure mode to watch \\
\midrule
Adaptive multiscale serialization & How should a model vary ordering granularity across dense and sparse regions? & Proteins mix compact motifs, flexible loops, and long-range contacts in the same structure & Over-optimizing local neighborhoods while obscuring distal functional relationships \\
Serialization-aware self-supervision & Which masking or reconstruction objectives preserve structural semantics after ordering? & Self-supervision is attractive when labels are scarce but structures are abundant & Learning order artifacts rather than biochemical regularities \\
Channel-rich structural tokens & Which atom, residue, or physicochemical channels should be bound to each serialized token? & Biomolecular tasks often depend on rich chemistry, not geometry alone & Token explosion that cancels the intended efficiency gain \\
Reversible interpretability & How much of the learned signal can be mapped back to 3D neighborhoods or residues? & Structural biologists need interpretable spatial hypotheses, not just scores & Losing traceability between token importance and physical structure \\
Hybrid sequence-structure systems & When should serialized 3D tokens interact with sequence-native foundation models? & Protein AI increasingly mixes sequence, structure, and function information & Weak fusion schemes that inherit the costs of both modalities without the benefits \\
\bottomrule
\end{tabular}
\caption{A publication-oriented agenda that turns the upgraded paper into a forward-looking contribution instead of a retrospective commentary.}
\label{tab:agenda}
\end{table*}


If this manuscript merely restated that the 2018 paper was prescient, it would not justify publication. The stronger contribution is to pair the reinterpretation with a measurable probe and then turn both into a forward agenda. Table~\ref{tab:agenda} summarizes the most defensible directions.

The first direction is \emph{adaptive multiscale serialization}. Modern systems already show that a single fixed ordering is rarely enough when local and global context both matter \citep{li2025noksr}. Biomolecular structures intensify this problem because the same protein may contain rigid cores, flexible loops, binding pockets, and long-range contacts. A publishable next-generation system should allow ordering granularity to vary with structural regime instead of freezing a single traversal for the whole object.

The second direction is \emph{serialization-aware self-supervision}. Models such as AtoMAE demonstrate that self-supervised structural pretraining is viable in biomolecular settings \citep{park2025atomae}. The open question is how masking and reconstruction objectives interact with the chosen ordering. A masking policy that respects structural neighborhoods may encourage different abstractions than one that respects only token position.

The third direction is \emph{channel-rich tokens}. One underappreciated strength of the 2018 framing was its focus on channel capacity. Future biomolecular systems should treat serialization not as geometry-only tokenization, but as the packaging of geometry together with chemistry. This is where representation design can still produce unique value even when backbones are commoditized.

The fourth direction is \emph{reversible interpretability}. If a model highlights a set of serialized tokens, how cleanly can those signals be projected back onto residues, atoms, or structural motifs? This matters for scientific credibility. Serialization is most attractive in science when it does not sever the path from model behavior back to physical explanation.

The fifth direction is \emph{hybrid sequence-structure modeling}. Sequence-native foundation models now dominate many biomolecular tasks, but they do not eliminate the need for explicit structure. Serialization offers one path toward a shared token space in which structural evidence can interact with sequence evidence without forcing both into the same raw representation. That possibility gives the 2018 idea a much broader modern relevance than its original experiments suggested.

Taken together, these directions define a publication path that is stronger than cosmetic revision. The manuscript is no longer asking readers to update an old benchmark table. It is offering a framework that can guide what the next empirical papers should optimize for, and an executable diagnostic that makes part of that framework measurable today.
